\documentclass[11pt,a4paper]{article}
\usepackage[UTF8]{ctex}
\usepackage[margin=2.5cm]{geometry}
\usepackage[numbers,sort&compress]{natbib}
\usepackage[colorlinks=true,citecolor=blue,linkcolor=blue,urlcolor=blue]{hyperref}
\usepackage{amsmath,amssymb,booktabs,longtable}
\title{弱网环境下跨境物流多模态感知与边缘协同研究进展}
\author{}
\date{}

\begin{document}
\maketitle

\begin{abstract}
跨境物流状态监测同时受到感知异构、网络间歇和端侧资源受限三类约束。现有工作分别在多模态异常检测、任务导向通信、分裂计算、时延容忍网络和不确定性预警方面取得进展，但大多默认模态同步、网络稳定或任务单一，尚未形成适合口岸—运输—仓储全过程的统一方法。本文梳理上述研究脉络，指出关键矛盾并非简单提高识别精度，而是使有限采样、传输与计算资源优先服务于异常证据保真，在模态缺失和场景漂移下维持告警可靠性。未来研究需建立状态可辨识、任务价值可度量、资源决策可约束和风险输出可校准的一体化框架，并通过跨路段、跨仓储条件和可重复弱网扰动进行验证。
\end{abstract}

\section{引言}

口岸查验、长距离运输和边境仓储构成连续而非孤立的物流过程。传感器能够记录振动、倾斜、温湿度、封志和门禁状态，图像与视频则提供包装破损、变形、烟火和入侵等直观证据。港口数字孪生研究把态势感知、数据分析和多方协同列为基础能力\cite{digitaltwinsfor2023}；物流网络研究也开始将异构物联网数据纳入多模态决策框架\cite{amultimodaldeep2025}。然而，现有方案通常把“采集—上传—分析”看作稳定链路，难以覆盖中亚跨境运输中的 2G/3G、长时延、频繁断连和端侧电量受限条件。

从算法角度看，跨境物流不是单一的图像异常检测问题。无线传感器异常融合工作说明，多源信号可以互补，但其收益依赖数据质量、时间对齐和融合机制\cite{multimodalfusionanomaly2024}。从系统角度看，它也不是普通的数据压缩问题：若压缩或丢弃的是异常发生前后的稀有片段，即使平均重建误差较低，也会直接损害告警。因而需要把感知表征、传输选择、边缘推理与风险输出放在同一任务目标下分析。

\section{物流多源感知的对象、证据与场景边界}

运输货物的异常证据具有多尺度特征。振动和倾斜反映冲击与姿态变化，电子封志和光感指示箱门开启，包装图像描述破损和变形；仓储侧的温湿度、烟雾、门禁与视频则共同描述环境风险。同一事件在不同模态上的响应并不同步，例如冲击先出现在惯性信号中，包装形变稍后才可由图像观察。若直接按固定时间戳拼接，可能把不同事件片段误认为同一状态。

已有工业过程数据集开始同时记录多源观测，为研究多模态过程异常提供了可复现实验基础\cite{amultimodaldataset2025}。工业控制系统异常数据集强调了物理过程与网络状态的联合监测\cite{industrialcontrolsystem2024}，真实二维/三维工业数据进一步说明视角、表面和结构异常不能由单一表征完全覆盖\cite{realiadd32025}。真实扰动基准还显示，训练分布上的高精度并不等价于面对噪声、遮挡和设备变化时的稳健性\cite{robustada2025}。这些结论可迁移到物流监测，但运输货物与固定产线仍有差别：设备安装姿态、道路工况和网络质量随路段变化，异常样本更稀缺，标注也更依赖事件日志。

因此，数据建模应保留“对象—事件—场景”三层语义：对象层描述货物、车辆和仓位；事件层对齐冲击、倾覆、非法开启、温湿度越界和烟火等证据；场景层记录路段、口岸、仓库与网络状态。该结构既能避免把网络丢包误判为传感异常，也为后续按事件价值选择上传内容提供依据。

\section{异步、缺失与漂移条件下的多模态异常表征}

异步多变量时序研究指出，观测间隔本身包含信息，简单插值可能掩盖异常边界；面向异步序列的检测与定位方法因此显式建模变量关系和观测时间\cite{practicalapproachto2021}。近期工作继续探索异步变量关系感知，使检测器在不同采样频率下保留跨变量依赖\cite{aaadasynchronousinter2025}。对跨境物流而言，这意味着高频振动、低频温湿度、事件触发封志和按需图像不能被强制压到同一采样率，而应通过时间编码、事件窗口和可信掩码进入统一表征。

模态缺失是另一常态。分布一致的模态恢复方法尝试从已观测模态推断缺失模态，同时约束恢复分布与完整数据一致\cite{distributionconsistentmodal2023}；鲁棒多传感器融合框架则通过训练期随机缺失提高测试期容错能力\cite{amultimodalsensor2022}。这类思路提示：恢复结果应作为带不确定性的辅助证据，而不是伪造“完整观测”。当封志或摄像头离线时，模型应降低相应证据权重，并输出缺失原因和置信区间。

工业异常检测中的注意力自编码与生成式融合可学习多模态正常模式\cite{mfganmultimodalfusion2024}，去噪型多模态方法进一步针对传感污染和三维噪声提高鲁棒性\cite{m3dmnrrgb2024}。实体—度量联合学习则从变量实体关系和距离结构两方面刻画异常\cite{madmelcombining2024}。这些方法若用于物流场景，需要额外处理跨车辆、跨仓库的分布偏移。道路传感异常研究表明，增量漂移检测有助于适应路况变化\cite{dynamicroadanomaly2024}；面向时间序列的 Transformer 漂移适应也说明，持续更新必须与异常记忆保护并行，否则模型可能把新近异常吸收为正常状态\cite{conceptdriftadaptation2022}。

综合来看，稳健表征应满足三项要求：一是对异步观测保持时间可辨识，二是对缺失模态输出显式置信度，三是把场景漂移与真实异常区分开。相应实验不应只报告平均 F1，还应设置随机缺失、连续断传、传感器偏置和跨路段迁移，并比较检测性能随扰动强度的退化曲线。

\section{面向任务价值的弱网通信与端边云协同}

传统传输目标倾向于最小化比特误码或信号重建失真，但物流告警关心的是异常判断是否改变。任务导向通信以推理损失而非像素/波形重建作为优化目标，机器到机器自编码通信已展示端到端联合编码的可行性\cite{anautoencoderbased2024}。基于语义三元组的方案把任务相关关系压缩为结构化信息\cite{taskorientedsemantic2023}，基础模型研究则讨论了语义先验带来的压缩能力与泛化风险\cite{taskorientedsemantic2024}。在物联网 DNN 推理中，目标导向通信直接以分类或检测性能约束传输\cite{goalorientedsemantic2024}；多模态语义中继研究进一步把不同模态的传输和边缘推理联合起来\cite{distributedtaskoriented2024}。

这些工作为弱网物流提供了“事件价值”视角：数据包的优先级不应只由生成时间决定，还应由其对异常后验概率、证据新颖性和告警时限的贡献决定。正常且重复的温湿度片段可以低频摘要，冲击前后的原始窗口、封志变化和包装图像则应高优先级缓存与重传。需要注意，任务导向压缩可能在训练分布外遗漏新型异常，因此必须保留最低限度的原始证据通道，并用保守阈值约束压缩强度。

端边云协同决定“在哪里算”。真实目标检测中，本地、边缘和分裂执行可根据上下文切换\cite{furciferacontext2024}；分裂计算结合早退机制，可在样本容易时端侧完成推理、困难样本再上传中间特征\cite{ispliteeimage2024}。边云协同目标检测把轻量端侧模型与卸载决策联合优化\cite{edgecloudcollaborative2024}，动态分裂框架则根据多任务负载调整切分点\cite{dynamicsplitcomputing2025}。车载边缘研究进一步将移动性与资源分配纳入卸载策略\cite{vehicularcomputingpower2025}，容错卸载方法提示切换策略必须考虑边缘节点不可用\cite{faulttolerantmobile2025}。

分裂联邦学习为多车辆/多仓库协作提供了不集中原始数据的路径，但通信轮次和非独立同分布数据会带来额外代价\cite{splitfederatedlearning2024}。人工智能定义的无线资源管理可联合信道、算力与卸载决策\cite{artificialintelligencedefined2024}。对本项目更合适的目标不是追求单一最优切分点，而是在带宽、能量、缓存和告警时限约束下，动态选择“端侧直接告警、上传特征复核、缓存原始证据等待补传”三种动作，并对策略失效设置安全回退。

\section{断续连接、低功耗与可靠交付}

时延容忍网络通过存储—携带—转发应对缺乏稳定端到端路径的环境。近期面向偏远区域的建模与设计说明，接触机会和缓存策略是交付率与时延的核心变量\cite{delaytolerantnetworks2025}；相关综述表明，复制数、路由知识和节点资源之间长期存在权衡\cite{reviewonrecent2022}。上下文感知转发利用节点状态提高选择性\cite{acontextaware2022}，沿海巡检等移动场景的路由工作则显示轨迹预测可以改善断续网络中的转发决策\cite{udtnrsa2023}。车载时延容忍网络研究还提醒，非合作节点或设备失效会破坏可靠交付假设\cite{emsefficientmonitoring2022}。

物流监测中的离线缓存和断点续传不能只保证文件最终到达，还要保证事件片段完整、顺序可恢复并可验证。缓存策略应优先保护异常前后窗口，重传单位应带有事件标识、时间范围、校验码和版本号；在长期断连时，端侧必须保留最低告警能力。低功耗方面，资源受限设备的能量感知自适应采样说明，采样率可随剩余能量与信号变化调整\cite{energyawareadaptive2023}；环境监测中的参数化自适应采样进一步证明，变化率和预测误差可驱动数据采集频率\cite{parametricmachinelearning2024}。

端侧检测同样需要预算约束。绿色物联网边缘异常检测把能耗纳入模型设计\cite{deeplearningdriven2024}，面向受限设备的离散化孤立森林则从内存和计算量降低无监督检测开销\cite{discretizedisolationforest2025}。分层工业物联网方案利用 MEC 承担复杂推理\cite{anomalydetectionfor2025}。可行的体系应形成两级闭环：端侧用轻量规则或模型完成即时筛查，边缘节点进行多模态复核，云端负责跨场景更新；任何一级离线时，上一层都能降级运行。

\section{仓储异常预警、可信输出与评价协议}

仓储风险往往由缓慢环境变化与突发事件共同构成。温湿度漂移可能先于质量风险，烟雾、门禁和视频异常则要求快速响应。异常与变点联合检测有助于区分渐变过程与瞬时异常\cite{anomalyandchange2023}，不确定性感知检测使模型能够把“高风险”和“证据不足”分开表达\cite{uncertaintyawaretime2024}。利用预测不确定性进行早期漂移检测，可在性能明显下降前触发复核\cite{earlyconceptdrift2025}；在线集成和漂移适应方法则提供持续更新机制\cite{ensemblebasedunsupervised2024}。近期无监督漂移适应工作继续探索非平稳多变量序列中的稳定更新\cite{unsupervisedconceptdrift2025}，鲁棒自适应工业物联网检测表明，多节点协作还需控制数据异质性\cite{fedramfederated2025}。

预警输出应至少包括异常类型、风险等级、置信度、证据来源和建议响应时限。固定阈值无法覆盖不同仓库和季节，宜通过校准集获得分场景阈值，并用期望校准误差、误报率和平均提前量共同评价。对包装图像和视频，可借鉴真实工业异常基准的扰动设置；对传感时序，可使用工业控制数据和多模态过程数据检验检测与定位能力\cite{industrialcontrolsystem2024,amultimodaldataset2025}。评价时还需报告上行字节量、包交付率、告警端到端时延、端侧能耗和峰值内存，避免算法精度与系统代价分开优化。

\section{讨论}

现有研究提供了可组合的技术部件，但直接拼接仍有四个风险。第一，多模态恢复可能生成看似合理却错误的证据，必须把恢复值与真实观测分离。第二，任务导向压缩依赖训练任务，未知异常可能被视为无关信息，需保留原始证据兜底。第三，卸载策略常以平均时延为目标，而跨境物流更关心告警截止时间和断连尾部风险。第四，在线适应可能产生灾难性遗忘，使持续出现的异常被错误归入正常模式。

此外，现有基准与真实跨境物流之间仍有明显鸿沟。多模态工业异常工作多聚焦固定设备，物流对象却持续移动；道路异常虽包含移动传感器，却不涉及货物多模态证据；港口与物流网络研究强调系统集成，但较少给出弱网下可证伪的算法假设。因此，需要构建可控弱网实验平台：重放真实或仿真的带宽、时延、丢包和断连轨迹，同时注入模态缺失、传感偏置和事件稀缺，使用相同事件流比较全量上传、固定压缩、规则触发和任务价值驱动方案。

\section{展望}

后续研究可沿三条相互约束的主线推进。其一，建立带可信掩码和时间编码的异步多模态状态空间，使缺失、延迟与异常在表示层可区分。其二，定义异常证据的边际任务价值，把采样、压缩、缓存、上传和模型切分写入统一受约束决策问题，并设置可靠交付与最低端侧告警能力。其三，构建运输—仓储跨场景风险模型，通过不确定性校准和漂移检测控制误报，在新的路段、季节和设备上验证泛化。若三条主线能够在共同的事件级评价协议下闭环，跨境物流监测才可能从“能采、能传、能识别”推进到“弱网下仍可信预警”。

\section{结论}

弱网跨境物流的核心科学挑战，是在异步、缺失和资源约束共同作用下保持异常证据可辨识与告警决策可靠。多模态鲁棒学习、任务导向通信、边云协同、时延容忍传输和可信预警分别提供了解题基础，但仍需由事件价值和约束优化贯通。面向新疆口岸、中亚运输和边境仓储的研究应坚持对象与场景边界，以跨场景退化、资源消耗和预警可靠性检验方法，而不以单一数据集精度代替系统有效性。

\bibliographystyle{gbt7714-nsfc}
\bibliography{跨境物流弱网多模态边缘协同_参考文献}

\end{document}
